\documentclass[12pt]{article}
\usepackage[margin=1in]{geometry}
\usepackage{amsmath}
\usepackage{hyperref}
\usepackage{parskip}
\usepackage{setspace}
\onehalfspacing

\title{%
  \textbf{Compartmental Epidemiological Models with\\
  Infection-Driven Viral Mutation}
}
\date{}

\begin{document}

\maketitle

\section*{Background and Motivation}

The classical SIR (Susceptible--Infected--Recovered) model and its variants
have proven surprisingly effective at describing the spread of infectious
diseases at the population level. A central assumption in the basic SIR
framework is that recovered individuals acquire permanent immunity, so that
each individual is infected at most once. In practice, two distinct mechanisms
can invalidate this assumption: (1) waning of immune memory over time, and (2)
viral mutation resulting in immune escape, whereby the pathogen evolves to
circumvent the defences of previously infected hosts.

The first mechanism---immunity waning---is well represented in the literature
through SIRS-type models, in which recovered individuals return to the
susceptible class at some fixed rate \cite{pease1987, casagrandi2006}.
The second mechanism---mutation-driven immune escape---has received considerably
less attention as a modelling target in its own right.
Of particular importance is the observation that mutations do not occur at a
fixed rate in time; rather, mutation opportunities arise proportionally to the
number of active infections, since each replication event inside a host is a
potential mutation event \cite{barrat2022, boni2006}.
This coupling between infection prevalence and mutation supply is fundamental:
it implies that the evolutionary trajectory of a pathogen is not independent
of epidemic dynamics, but is instead shaped by them.

The COVID-19 pandemic has made this interaction especially visible.
Successive variants of concern---Alpha, Delta, Omicron, and its
subvariants---demonstrated that immune escape, rather than waning immunity
alone, was the dominant driver of renewed epidemic waves \cite{naturerev2023, covid_past2022}.
The pandemic also generated an unprecedented volume of genomic sequence data,
providing new opportunities to calibrate and validate mathematical models of
viral evolution at the population level \cite{ngonghala2025}.

\section*{Project Goals}

This project will explore several approaches to incorporating
infection-driven viral mutation into compartmental epidemiological models.
The central modelling principle is that the rate at which the virus accumulates
immune-escape mutations is proportional to the current infected population
$I(t)$, rather than being a fixed external parameter.
This is in contrast to the standard SIRS phenomenology, which implicitly treats
antigenic drift as a clock running independently of epidemic dynamics
\cite{pease1987, yaari2014}.

Specific objectives include:

\begin{enumerate}
  \item \textbf{Model formulation.} Develop and compare two or three
    ODE-based compartmental models of increasing complexity. Candidate
    architectures include:
    \begin{itemize}
      \item A modified SIRS model in which the rate of transition from $R$
        to $S$ is proportional to $I(t)$, extending the Pease framework
        \cite{pease1987} with endogenous mutation supply.
      \item A SIRC-type model \cite{casagrandi2006, li2017} augmented with
        an infection-dependent cross-immunity erosion rate.
      \item A combined ODE-PDE model in which immunity of the recovered population varies in time via an advection equation.
%      \item A two-strain SIR model \cite{an2020} in which a mutant strain
%        is generated at a rate proportional to the prevalence of the
%        wild-type strain, following the multiscale coupling approach of
%        recent work \cite{multiscale2026}.
    \end{itemize}

  \item \textbf{Equilibrium and stability analysis.} For each model,
    compute disease-free and endemic equilibria, determine the basic
    reproduction number $R_0$, and analyse local (and where possible global)
    stability. A key question is whether the system admits a critical
    population threshold below which epidemics die out and above which they
    become cyclic or endemic.

  \item \textbf{Bifurcation analysis.} Identify bifurcations in the
    qualitative dynamics as parameters are varied---particularly the
    infection rate, mutation rate, and cross-immunity---using both analytical
    and numerical methods. The SIRC model is already known to exhibit
    rich bifurcation structure including chaos \cite{casagrandi2006}, and
    similar or richer behaviour is anticipated when mutation rate is coupled
    to prevalence. Multi-strain models with cross-immunity are known to
    produce epidemic cycling \cite{zhang2016}.

  \item \textbf{Numerical simulation.} Solve the ODE systems numerically
    over biologically relevant parameter ranges, generating time series of
    $S(t)$, $I(t)$, $R(t)$ and any additional compartments. Classify the
    qualitative regimes: extinction, single epidemic, sustained endemic,
    or recurrent cyclic outbreaks.

  \item \textbf{Comparison with data (if time permits).} Compare model
    predictions with the succession of COVID-19 variant waves, using the
    hypothesis-testing framework of \cite{ngonghala2025} as a template.
\end{enumerate}

\section{Models}

\subsection{ODE-PDE model}

\begin{subequations}
\begin{align}
    x'(t) & = - \beta x y + v z(t,1) \\
    y'(t) & = \beta x y - \gamma y + \beta y \int_0^1 f(c) z(c,t) dc \\
    z_t + v z_s & = -\beta y f(c) z(c,t).
\end{align}
\end{subequations}
Here $c\in[0,1]$ is the cross-immunity.  For the \emph{antigenic drift velocity} $v$,
we could take a constant value or (in line with the idea of this project) relate it
to the rate of new infections:
\begin{align}
    v(t) & = \delta \beta x y,
\end{align}
where $\delta$ is the mutation rate per infection.
For $z$ we require a boundary condition at $s=0$, which is given by the
influx of people leaving the infected pool:
\begin{align}
    z(t,c=0) & = \gamma y(t).
\end{align}
Meanwhile, $f(c)$ determines how beneficial cross-immunity is.  If we take $f(c)=0$,
then any value $c<1$ gives full immunity.  Alternatively, $f$ could be any increasing
function, going from a value of zero to one, such as $f(c)=c$, $f(c)=c^2$, etc.

\bibliographystyle{plain}
\bibliography{references}

\end{document}
